% Paper 10: The Crypto Fear Channel — Asymmetric BTC-Equity Volatility Spillover
% body_v1.tex: §1 Introduction (carry-over from body_v0_intro.tex 2026-04-20) +
%              §2 Literature Review (NEW, drafted 2026-04-28 main thread)
% Numbers all sourced from experiments/k1025/k1025_results.json.

\documentclass[11pt,a4paper]{article}
\usepackage{graphicx}
\usepackage{amsmath,amssymb}
\usepackage{natbib}
\usepackage[margin=1in]{geometry}
\usepackage{booktabs}
\usepackage{url}

\title{The Crypto Fear Channel:\\ Asymmetric, Tail-Concentrated, and Regime-Dependent Volatility Spillover from Bitcoin to Equity Markets}

\author{Yi-Hao Lai\thanks{Department of Finance, Da-Yeh University, Taiwan. Email: yihao.lai@gmail.com.}}

\date{\today}

\begin{document}
\maketitle

\begin{abstract}
\noindent Using daily returns on SPY, BTC-USD, and VIX from 2015-02 to 2026-04 ($N = 2{,}812$), we document three stylized facts about the crypto-to-equity volatility spillover. First, the spillover is \emph{asymmetric}: Bitcoin downside volatility Granger-causes VIX at lags 1--5, whereas Bitcoin upside volatility does not (symmetric Granger tests covering lags 1--10 corroborate the broader directional result).% source: experiments/k1025/k1025_results.json .asymmetric_granger (lags 1-5) + .granger_causality.btc_rv_to_vix (symmetric 1-10)
 Second, the spillover is \emph{tail-concentrated with sign reversal}: the quantile-regression coefficient of VIX on BTC realized variance is significantly \emph{negative} at lower quantiles ($-2.86$ at $\tau = 0.05$; $-2.34$ at $\tau = 0.25$), indicating bull-market decoupling, then turns positive at the median ($+2.61$) and amplifies to $+22.31$ at the 95th percentile, yielding an $8.5\times$ upper-tail amplification relative to the median. Third, the spillover is \emph{regime-dependent}: in a five-subperiod breakdown, Granger causality is statistically significant only during 2020 ($F = 11.05$, $p < 0.001$) and non-significant in 2015--2017, 2018--2019, 2021--2022, and 2023--2026. In the Diebold-Yilmaz framework, Bitcoin is a \emph{net receiver}, rehabilitating the narrative from ``crypto as a fear originator'' to ``crypto as a fear amplifier'' conditional on pre-existing equity stress. Out-of-sample forecasts, however, show that adding BTC realized volatility does not statistically improve an AR model of VIX (Diebold-Mariano $t = -0.98$, $p = 0.33$), highlighting that Granger causality is a necessary but not sufficient condition for forecastability. Taken together, our findings reframe the crypto-equity spillover as a \emph{crisis-time amplifier}: informative for ex-post risk attribution and margin design, but inadequate as an ex-ante predictor of equity fear. \\

\noindent\textbf{Keywords:} volatility spillover, Bitcoin, VIX, asymmetric causality, quantile regression, Diebold-Yilmaz, crypto-equity linkage \\
\textbf{JEL:} G11, G15, G17, C22
\end{abstract}

\section{Introduction}

Since the institutionalization of Bitcoin trading and the 2024 approval of U.S.\ spot BTC ETFs, the question of whether cryptocurrency markets transmit fear to traditional equity markets has acquired first-order practical importance for portfolio risk management, prudential supervision, and margin design. Yet the empirical literature has yielded an uneasy consensus. On one hand, correlation-based studies report rising BTC-equity co-movement in stress periods, suggesting that cryptocurrency is no longer a standalone asset class.\footnote{See, e.g., \citet{bouri2020}, \citet{corbet2018}, and \citet{matkovskyy2019} for the early post-2017 literature.} On the other hand, return-predictability tests often fail to find that Bitcoin-based signals improve forecasts of equity volatility at practical horizons. This tension motivates a central question: is the crypto-equity fear channel \emph{real but nonlinear}, \emph{real but regime-dependent}, or an artifact of tail-period correlation without predictive content?

We resolve this tension by decomposing the spillover along three dimensions that the prior literature has typically examined in isolation: directional asymmetry, tail dependence, and regime stability. Each of these dimensions yields a stylized fact that materially changes how one should interpret correlation-based evidence.

\paragraph{Asymmetry.} Using the asymmetric Granger-causality framework of \citet{hatemi2012}, we separate Bitcoin positive and negative realized volatility and test whether each Granger-causes VIX at lags 1 through 5.% source: experiments/k1025/k1025_results.json .asymmetric_granger.btc_neg_to_vix (lags 1-5)
 We find that the negative branch carries the entire spillover: Bitcoin downside volatility Granger-causes VIX at every lag tested (1--5), while the upside branch is uniformly non-significant across the same lag window. This is not merely a statistical artifact of a leverage effect --- the asymmetry is corroborated by standard symmetric Granger tests over the broader 1 through 10 lag range, where BTC-to-VIX causality is significant at lags 2--10.% source: experiments/k1025/k1025_results.json .granger_causality.btc_rv_to_vix (symmetric lags 1-10)
 Our interpretation is behavioral: retail-heavy crypto markets transmit \emph{fear}, not \emph{euphoria}, into the institutional equity complex.

\paragraph{Tail concentration with sign reversal.} Quantile regressions of VIX on BTC realized variance reveal a richer structure than a simple tail amplification: the slope \emph{reverses sign} across the VIX distribution. At the lower quantiles the coefficient is significantly \emph{negative} ($-2.86$ at $\tau = 0.05$; $-2.34$ at $\tau = 0.25$), consistent with a bull-market coexistence regime in which rising Bitcoin volatility coincides with \emph{compressing} equity fear. The slope turns positive at the median ($+2.61$ at $\tau = 0.5$) and amplifies sharply in the upper tail ($+22.31$ at $\tau = 0.95$), an approximately $8.5\times$ amplification of the upper tail relative to the median ($\tau = 0.95 / \tau = 0.5 \approx 8.54$).% source: experiments/k1025/k1025_results.json .quantile_regression (tau=0.05,0.25,0.5,0.95)
 The spillover is therefore not a uniform channel, nor even a simple ``flat-then-spikes'' tail phenomenon, but a regime with a \emph{directional switch} at the median of VIX: BTC volatility dampens equity fear in calm markets and amplifies it in panic markets. Conventional OLS or DCC-GARCH estimates --- which collapse this sign-reversing heterogeneity into a single slope --- systematically mischaracterize both the bull-market decoupling and the crisis-time amplification of crypto volatility on equity fear.

\paragraph{Regime dependence.} A five-subperiod breakdown reveals that Granger significance is concentrated entirely in 2020 ($F = 11.05$, $p < 10^{-6}$). Pre-mania (2015--2017), crypto-winter (2018--2019), bull-bear (2021--2022), and recovery-plus-ETF (2023--2026) all fail to reject non-causality. Combined with the finding from \citet{diebold2012} spillover-index analysis that Bitcoin is a \emph{net receiver} rather than net sender, this suggests that the spillover is not a permanent feature of cross-market architecture but a \emph{conditional amplification mechanism} that activates only during pre-existing equity stress.

\paragraph{The forecastability gap.} A natural next question is whether these three stylized facts translate into out-of-sample predictive power. They do not. Comparing an AR($p$) specification of VIX against one augmented with BTC realized volatility over an OOS window of 1,826 days (2019--2026), the Diebold-Mariano test returns $t = -0.98$, $p = 0.33$ --- statistically indistinguishable from zero, and below the \citet{harvey2016} $|t| > 3$ threshold for newly proposed factors. This result echoes a broader methodological point: \emph{Granger causality is a necessary but not sufficient condition for forecastability}. An informative signal that appears only in tail events and only during rare regimes may be too sparse to outperform a well-specified autoregressive benchmark on average over a long OOS sample.

We make three contributions to the literature. First, by combining asymmetric Granger causality, quantile regression for tail dependence, Diebold-Yilmaz spillover direction, and an honest out-of-sample forecasting test in a single framework, we show that each dimension tells a distinct story that prior single-technique studies have missed. Second, the regime-dependent finding --- that COVID-2020 is a structural watershed --- implies that volatility spillover from crypto is a \emph{crisis-time amplifier} rather than a steady-state channel, with direct implications for margin and capital requirements during regime transitions. Third, the predictive-power null, reported transparently rather than hidden, contributes to a growing methodological awareness in the volatility forecasting literature that in-sample significance and out-of-sample usefulness can diverge substantially \citep{harvey2016, conrad2020}.

The remainder of the paper is organized as follows. Section~\ref{sec:lit} surveys the crypto-equity spillover literature and positions our contribution. Section 3 describes the data and preliminary diagnostics. Section 4 details the five methodological building blocks. Sections 5 and 6 present the main results and robustness checks. Section 7 reports the out-of-sample forecasting evaluation. Section 8 discusses mechanisms and policy implications. Section 9 concludes.

\section{Literature Review}
\label{sec:lit}

We organize the related literature along three threads: (i) the empirical literature on crypto-equity volatility spillover and the safe-haven debate, (ii) the methodological literature underlying our four building blocks (asymmetric Granger, quantile regression, Diebold-Yilmaz, Diebold-Mariano), and (iii) the recent calls for honest out-of-sample evaluation in volatility forecasting. Each thread has progressed independently; our contribution is to combine them within a single framework that decomposes the spillover into asymmetry, tail concentration, regime conditionality, and predictive power.

\subsection{Crypto-equity volatility spillover and the safe-haven debate}

The empirical literature on Bitcoin's role in equity portfolios has evolved through three phases. The first phase, exemplified by \citet{corbet2018} and \citet{matkovskyy2019}, documented that cryptocurrency price movements respond meaningfully to macro-financial news and that contagion-style co-movement with equity markets emerges in stress periods, contradicting the early ``digital gold'' framing. The second phase, motivated by COVID-19 and reviewed in \citet{bouri2020}, \citet{conlon2020}, \citet{shahzad2019}, and \citet{klein2018}, sharpened the safe-haven question with mixed verdicts: wavelet- and DCC-based correlation analysis showed that Bitcoin's risk profile is closer to a high-volatility risky asset than to gold, and that its safe-haven properties are at best regime-dependent. The third phase, post-2020, has documented increasing institutional integration: \citet{yarovaya2022} report that information transmission between cryptocurrency and traditional asset markets intensified during the pandemic, and \citet{akyildirim2020} link implied-volatility shocks to cryptocurrency returns.

Three gaps motivate our analysis. First, most prior studies test \emph{symmetric} causality or \emph{symmetric} correlation, conflating the upside and downside volatility channels that we will show carry economically distinct information. Second, the tail dependence of the spillover is typically inferred indirectly from rolling DCC or wavelet coherence, rather than directly estimated through quantile regression. Third, sample windows often end at 2020 or include 2020 as a single observation, leaving the post-COVID regime-stability question unaddressed --- a gap that matters precisely because COVID-2020 will turn out to be the structural watershed of the spillover. Our analysis closes these gaps by replacing each symmetric or correlation-based summary with a direct, regime-aware, asymmetric counterpart, while extending the sample through April 2026.

\subsection{Methodological building blocks}

Our four building blocks each have a settled methodological literature. Asymmetric Granger causality, the technique behind our headline asymmetry result, was formalized by \citet{hatemi2012}, who decompose each series into positive and negative cumulative innovations and test their respective Granger predictability. Earlier nonparametric extensions of Granger causality \citep{diks2006} and asymmetric volatility-spillover tests \citep{hong2001} provide complementary diagnostic tools that we use as robustness checks. Each of these methods was developed to address the specific concern that linear Granger tests can mask sign-dependent transmission --- precisely the asymmetry the present paper documents in the BTC-to-VIX channel.

Quantile regression, our tool for tail dependence, originates with \citet{koenker1978} and has subsequently been extended into systemic-risk applications. Most relevant to our setting is \citet{adrian2016}'s CoVaR framework, in which the conditional value-at-risk of one institution given another's stress is estimated through quantile regression on returns. We adopt the quantile-regression framework rather than CoVaR \emph{per se} because our object of interest is the conditional response of equity \emph{volatility} (VIX), not the conditional value-at-risk of equity returns; but the underlying statistical engine is the same.

The Diebold-Yilmaz spillover index is now a canonical tool for measuring directional volatility transmission across asset classes. Its development spans \citet{diebold2009}, \citet{diebold2012}, and \citet{diebold2014network}, and a recent IMF policy note \citep{iyer2022} applies the framework specifically to cryptocurrency-equity spillovers. We follow the 2012 specification (variance decomposition of a generalized VAR) and report rolling 252-day spillover indices to assess directional connectedness over time.

Finally, our forecasting evaluation rests on the \citet{diebold1995} test of equal predictive accuracy and the \citet{harvey2016} multiple-testing threshold of $|t| > 3$ for newly proposed factors. The latter has become a methodological norm for newly proposed predictors in asset pricing, and we adopt it in the volatility-forecasting setting to avoid the inflated false-positive rate associated with conventional $|t| > 1.96$ inference under multiple testing.

\subsection{Honest out-of-sample evaluation in volatility forecasting}

A recurring tension in the volatility-forecasting literature is the gap between in-sample statistical significance and out-of-sample economic value. \citet{conrad2020}'s GARCH-MIDAS application demonstrates that incorporating macro-economic information through long-run components can deliver in-sample fit improvements that fail to translate into forecast accuracy gains. The cumulative evidence reviewed there motivates an evaluation discipline in which a new predictor must clear both (i) in-sample significance under the Harvey threshold and (ii) out-of-sample improvement under Diebold-Mariano, before being promoted to a structural claim. The present paper adopts this discipline and reports a transparent out-of-sample null --- BTC realized volatility fails to improve AR-VIX forecasts at $t = -0.98$ --- alongside its in-sample asymmetry, tail-amplification, and regime-stability findings. The combination is, to our knowledge, the first such honest joint reporting in the crypto-equity spillover literature.

% --- References (temporary stub; full bib to follow in main draft) ---
\begin{thebibliography}{99}

\bibitem[Adrian and Brunnermeier, 2016]{adrian2016}
Adrian, T. and Brunnermeier, M.K. (2016).
\newblock CoVaR.
\newblock {\em American Economic Review}, 106(7):1705--1741.
\newblock \url{https://doi.org/10.1257/aer.20120555}

\bibitem[Akyildirim et al., 2020]{akyildirim2020}
Akyildirim, E., Corbet, S., Lucey, B., Sensoy, A., and Yarovaya, L. (2020).
\newblock The relationship between implied volatility and cryptocurrency returns.
\newblock {\em Finance Research Letters}, 33:101212.
\newblock \url{https://doi.org/10.1016/j.frl.2019.06.010}

\bibitem[Bouri et al., 2020]{bouri2020}
Bouri, E., Shahzad, S.J.H., Roubaud, D., Kristoufek, L., and Lucey, B. (2020).
\newblock Bitcoin, gold, and commodities as safe havens for stocks: New insight through wavelet analysis.
\newblock {\em The Quarterly Review of Economics and Finance}, 77:156--164.
\newblock \url{https://doi.org/10.1016/j.qref.2020.03.004}

\bibitem[Conlon and McGee, 2020]{conlon2020}
Conlon, T. and McGee, R. (2020).
\newblock Safe haven or risky hazard? Bitcoin during the Covid-19 bear market.
\newblock {\em Finance Research Letters}, 35:101607.
\newblock \url{https://doi.org/10.1016/j.frl.2020.101607}

\bibitem[Conrad and Kleen, 2020]{conrad2020}
Conrad, C. and Kleen, O. (2020).
\newblock Two are better than one: Volatility forecasting using multiplicative component GARCH-MIDAS models.
\newblock {\em Journal of Applied Econometrics}, 35(1):19--45.

\bibitem[Corbet et al., 2018]{corbet2018}
Corbet, S., Larkin, C., Lucey, B., Meegan, A., and Yarovaya, L. (2018).
\newblock Cryptocurrency reaction to FOMC announcements.
\newblock {\em Economics Letters}, 165:28--34.
\newblock \url{https://doi.org/10.1016/j.econlet.2018.01.004}

\bibitem[Diebold and Mariano, 1995]{diebold1995}
Diebold, F.X. and Mariano, R.S. (1995).
\newblock Comparing predictive accuracy.
\newblock {\em Journal of Business \& Economic Statistics}, 13(3):253--263.

\bibitem[Diebold and Yilmaz, 2009]{diebold2009}
Diebold, F.X. and Yilmaz, K. (2009).
\newblock Measuring financial asset return and volatility spillovers, with application to global equity markets.
\newblock {\em The Economic Journal}, 119(534):158--171.
\newblock \url{https://doi.org/10.1111/j.1468-0297.2008.02208.x}

\bibitem[Diebold and Yilmaz, 2012]{diebold2012}
Diebold, F.X. and Yilmaz, K. (2012).
\newblock Better to give than to receive: Predictive directional measurement of volatility spillovers.
\newblock {\em International Journal of Forecasting}, 28(1):57--66.
\newblock \url{https://doi.org/10.1016/j.ijforecast.2011.02.006}

\bibitem[Diebold and Yilmaz, 2014]{diebold2014network}
Diebold, F.X. and Yilmaz, K. (2014).
\newblock On the network topology of variance decompositions: Measuring the connectedness of financial firms.
\newblock {\em Journal of Econometrics}, 182(1):119--134.
\newblock \url{https://doi.org/10.1016/j.jeconom.2014.04.012}

\bibitem[Diks and Panchenko, 2006]{diks2006}
Diks, C. and Panchenko, V. (2006).
\newblock A new statistic and practical guidelines for nonparametric Granger causality testing.
\newblock {\em Journal of Economic Dynamics and Control}, 30(9--10):1647--1669.
\newblock \url{https://doi.org/10.1016/j.jedc.2005.08.008}

\bibitem[Harvey et al., 2016]{harvey2016}
Harvey, C.R., Liu, Y., and Zhu, H. (2016).
\newblock \ldots and the cross-section of expected returns.
\newblock {\em Review of Financial Studies}, 29(1):5--68.

\bibitem[Hatemi-J, 2012]{hatemi2012}
Hatemi-J, A. (2012).
\newblock Asymmetric causality tests with an application.
\newblock {\em Empirical Economics}, 43(1):447--456.
\newblock \url{https://doi.org/10.1007/s00181-011-0484-x}

\bibitem[Hong, 2001]{hong2001}
Hong, Y. (2001).
\newblock A test for volatility spillover with application to exchange rates.
\newblock {\em Journal of Econometrics}, 103(1--2):183--224.
\newblock \url{https://doi.org/10.1016/S0304-4076(01)00043-4}

\bibitem[Iyer, 2022]{iyer2022}
Iyer, T. (2022).
\newblock Cryptic connections: Spillovers between crypto and equity markets.
\newblock {\em IMF Global Financial Stability Notes}, 2022/01.
\newblock \url{https://doi.org/10.5089/9781616358068.065}

\bibitem[Klein et al., 2018]{klein2018}
Klein, T., Pham Thu, H., and Walther, T. (2018).
\newblock Bitcoin is not the new gold --- A comparison of volatility, correlation, and portfolio performance.
\newblock {\em International Review of Financial Analysis}, 59:105--116.
\newblock \url{https://doi.org/10.1016/j.irfa.2018.07.010}

\bibitem[Koenker and Bassett, 1978]{koenker1978}
Koenker, R. and Bassett, G. (1978).
\newblock Regression quantiles.
\newblock {\em Econometrica}, 46(1):33--50.
\newblock \url{https://doi.org/10.2307/1913643}

\bibitem[Matkovskyy and Jalan, 2019]{matkovskyy2019}
Matkovskyy, R. and Jalan, A. (2019).
\newblock From financial markets to Bitcoin markets: A fresh look at the contagion effect.
\newblock {\em Finance Research Letters}, 31:93--97.
\newblock \url{https://doi.org/10.1016/j.frl.2019.04.007}

\bibitem[Shahzad et al., 2019]{shahzad2019}
Shahzad, S.J.H., Bouri, E., Roubaud, D., Kristoufek, L., and Lucey, B. (2019).
\newblock Is Bitcoin a better safe-haven investment than gold and commodities?
\newblock {\em International Review of Financial Analysis}, 63:322--330.
\newblock \url{https://doi.org/10.1016/j.irfa.2019.01.002}

\bibitem[Yarovaya et al., 2022]{yarovaya2022}
Yarovaya, L., Brzeszczy\'nski, J., Goodell, J.W., Lucey, B., and Lau, C.K.M. (2022).
\newblock Rethinking financial contagion: Information transmission mechanism during the COVID-19 pandemic.
\newblock {\em Journal of International Financial Markets, Institutions and Money}, 79:101589.
\newblock \url{https://doi.org/10.1016/j.intfin.2022.101589}

\end{thebibliography}

\end{document}
